\chapter{Conclusions}
\label{ch:conclusions}

This Thesis asked how a quantum sensor can be assessed over the regimes in which it must operate.  QFI gives the best sensitivity allowed by a parameterized quantum state, but a practical design must also account for state preparation, restricted access, noise, measurement, calibration, and the data that are actually recorded.

\section{Main Findings}

The theoretical part implemented a mixed-state QFI pipeline based on TQFI and SSQFI bounds.  VQSE supplies the dominant state components needed by TQFI, while VQFIE uses fidelity measurements to estimate and optimize metrological bounds.  These methods reduce the need for full classical diagonalization and offer a route for studying larger mixed probes.

The TFIM calculations treated the transverse field as the signal of a finite quantum magnetometer.  Restricted access could not increase QFI within the same preparation family, and keeping more of the probe recovered more of the full-system response.  In the tested cases, keeping all but one qubit recovered essentially all of the full-system peak.  Fixed depolarization lowered local QFI monotonically in the tested range.  The thermal results showed that the best preparation depends on the expected operating regime.  A finite-temperature Gibbs preparation can outperform the reduced ground-state preparation over broad field intervals away from the ground-state optimum.  At $\beta=5$ in the studied $N=6$, $n=4$ periodic system, its peak local QFI was also $23.5\%$ higher.  This advantage comes from the field-dependent populations of low-energy states.  It compares distinct preparation families and does not mean that a fixed thermal channel increases QFI.

The optical part reached a similar design lesson.  Direct imaging loses sensitivity in the high-contrast sub-Rayleigh regime, while SPADE places useful separation information in low-order spatial modes.  In the theoretical calculation per surviving photon, full ideal SPADE kept at least $0.99899$ of the source-state QFI over the tested regular settings.  The benefit can fall sharply after binary aggregation and readout confusion, which shows why the measurement and detector layers must be included in the analysis.

The experimental SPADE study also changed the methodological direction of the Thesis.  The initial question was if the completed record could support a TQFI analysis, but the archived aggregate counts support neither variational diagonalization nor reconstruction of the reduced optical state.  This is a limit of the archived record rather than a fundamental limit of SPADE or TQFI.  A model-based TQFI can be calculated after specifying the source and every channel, but it is then conditional on those choices and is not reconstructed from the experiment.

The main experimental contribution was therefore developed in two layers.  First, a classical calibration model was fitted to the population-level setting means.  It predicted held-out $100$-repeat means when checkerboard cells, complete distance rows, and complete source-ratio columns were excluded from fitting.  Under the stated calibration and readout conventions, the identifiable scale is the conditional throughput $T=1304.25$ per $10\,\mathrm{ms}$, with a grouped cross-fit envelope of $[1251.49,1359.47]$, rather than an absolute efficiency.

Second, the complete experiment was organized as a sequential quantum-to-classical model.  The framework separates source preparation, optical channels, accessible modes, SPADE measurement, classical readout, calibration, and count statistics.  This separation preserves probability, distinguishes coherence effects from population effects, and provides checkpoints for conditional QFI and CFI calculations.  It also makes parameter provenance and non-identifiability explicit.  The fitted data constrain the observable classical projection of this sequence.  They do not fit or reconstruct the complete quantum channels.

\section{Overall Lesson and Outlook}

Across both applications, the central lesson is that useful sensitivity is a property of the full sensing procedure and its operating range.  Mixedness can arise through different physical processes with different QFI.  Access to only part of a probe can hide information.  A nearly optimal theoretical measurement can lose its advantage through coarse readout, uncertain calibration, or incomplete records.  The value of the channel model is therefore not that it introduces more parameters.  Its value is that it connects each loss of information to a physical or statistical stage and shows which conclusions the available observations can support.

For magnetometry, the next steps are larger-system tests, hardware implementations of the variational bounds, and preparation studies matched to a known field range rather than only to the largest QFI peak.  For SPADE, a new experiment could be planned with ASI Matera to record separate HG outputs, incident flux, background, failure outcomes, and phase-sensitive observables.  These additions would first allow the CFI of the actual receiver to be measured.  Measurements in additional modal bases and independent channel calibrations would then provide the state and channel information needed for a sound experimental comparison with QFI and TQFI.  As a secondary future direction, the proposed VQ-SPADE scheme could add a programmable modal rotation and use resolved SPADE outputs as an optical VQSE readout.  It was not implemented or tested in the present work.  The framework developed here therefore serves both as an interpretation of the current data and as a design guide for the next experiment.
